\chapter{Epiglottitis}
\index{Epiglottitis}

\section*{Definition and Overview}
Epiglottitis is a serious, life-threatening infection and swelling of the epiglottis, the flap of cartilage at the base of the tongue that covers the windpipe when swallowing. When the epiglottis swells, it can block the airway and cause a sudden inability to breathe. Epiglottitis is most commonly caused by the bacterium Haemophilus influenzae type b (Hib), and it was once a common illness of young children, but routine vaccination has made it rare. It can also affect adults, and it requires immediate emergency treatment to secure the airway.

\section*{Epidemiology}
Since the introduction of the Hib vaccine in the 1990s, epiglottitis in children has become rare in many countries. Cases still occur in unvaccinated children and in adults, in whom the infection may be caused by other bacteria, such as streptococci and staphylococci. It is a medical emergency at any age, and prompt recognition and treatment are essential because the airway can close rapidly.

\section*{Aetiology and Pathophysiology}
Epiglottitis is caused by infection of the epiglottis and surrounding tissues.
\begin{itemize}
    \item \textbf{Haemophilus influenzae type b:} the classic cause in children, prevented by vaccination
    \item \textbf{Other bacteria:} group A streptococcus, Streptococcus pneumoniae, and Staphylococcus aureus in adults
    \item \textbf{Transmission:} the bacteria spread by droplets from the nose and throat
    \item \textbf{Mechanism:} infection causes rapid, severe swelling of the epiglottis and the tissues around the voice box
    \item \textbf{Airway obstruction:} the swollen epiglottis blocks the entrance to the windpipe, making breathing difficult and then impossible
    \item \textbf{Onset:} symptoms can progress from a sore throat to airway obstruction within hours
\end{itemize}

\section*{Clinical Features}
Epiglottitis has a distinctive presentation.
\begin{itemize}
    \item Sudden onset of a high fever and severe sore throat
    \item Difficulty and pain on swallowing, with drooling
    \item Difficulty breathing, with a muffled or "hot potato" voice
    \item The child sits forward, with the neck extended and mouth open, trying to breathe
    \item No cough, which helps distinguish it from croup
    \item Restlessness, anxiety, and a blue tinge as oxygen falls
    \item In adults: severe sore throat and difficulty swallowing, progressing to breathing difficulty
\end{itemize}

\section*{Differential Diagnosis}
Other conditions cause breathing difficulty in children and adults.
\begin{itemize}
    \item Croup, a viral infection of the airway with a barking cough
    \item Severe tonsillitis and peritonsillar abscess
    \item A foreign body in the airway
    \item Severe allergic reactions causing airway swelling
    \item Asthma and other causes of breathing difficulty
    \item Bacterial tracheitis
\end{itemize}

\section*{When to Seek Medical Care}
Epiglottitis is an emergency. Any child or adult with sudden difficulty swallowing, drooling, muffled voice, or difficulty breathing must be treated as a medical emergency.

\textbf{Contact your local emergency services for:}
\begin{itemize}
    \item Difficulty breathing, stridor, or a blue tinge
    \item Drooling and an inability to swallow
    \item A muffled voice with a high fever
    \item A child sitting forward, unable to breathe comfortably
    \item Collapse or loss of consciousness
\end{itemize}

\section*{Diagnosis and Investigations}
Diagnosis is made carefully, without disturbing the child's airway.
\begin{itemize}
    \item \textbf{Clinical recognition:} the classic features are recognised, and the airway is protected
    \item \textbf{No throat examination in the clinic:} examining the throat can trigger complete airway obstruction
    \item \textbf{Hospital assessment:} in a controlled setting, doctors may examine the epiglottis
    \item \textbf{Blood tests:} after the airway is secure
    \item \textbf{Imaging:} X-rays may be used in stable patients to support the diagnosis
\end{itemize}

\section*{Management}
Treatment happens in an emergency department with an intensive care team.
\begin{itemize}
    \item \textbf{Protect the airway:} the child is kept calm and comfortable, often in the parent's arms, and is not examined or disturbed
    \item \textbf{Securing the airway:} a breathing tube is placed by an experienced specialist, often in the operating theatre
    \item \textbf{Antibiotics:} intravenous antibiotics against the likely bacteria
    \item \textbf{Corticosteroids:} reduce swelling
    \item \textbf{Intensive care:} the child is monitored closely until the swelling settles, usually a few days
    \item \textbf{Recovery:} the breathing tube is removed when the epiglottis is no longer swollen
\end{itemize}

\section*{Complications}
\begin{itemize}
    \item Complete airway obstruction and death, if not treated promptly
    \item Spread of infection, including pneumonia, meningitis, and blood poisoning
    \item Complications of intensive care
    \item Rarely, permanent airway damage
\end{itemize}

\section*{Prognosis and Outcomes}
With prompt emergency treatment, the outlook for epiglottitis is very good, and most people recover completely within about a week. The disease is preventable by vaccination, and cases in many countries have fallen dramatically since the Hib vaccine was introduced. The key to a good outcome is recognising the emergency and seeking help immediately.

\section*{Prevention and Self-Care}
\begin{itemize}
    \item Keep childhood immunisations up to date, including the Hib vaccine
    \item Recognise the warning signs and call for help immediately
    \item Do not try to examine the throat or give food or drink
    \item Keep the person calm and in a comfortable position while waiting for help
    \item Adults should be aware that epiglottitis can also occur in adulthood
\end{itemize}

\section*{Sources and Further Reading}
The following reputable sources provide further information for healthcare students and patients seeking reliable, current advice about epiglottitis.
\begin{itemize}
    \item The Royal Children's Hospital Melbourne. \textit{Epiglottitis fact sheet.} https://www.rch.org.au/
    \item Healthdirect Australia. \textit{Epiglottitis.} https://www.healthdirect.gov.au/epiglottitis
    \item NSW Health. \textit{Epiglottitis information.} https://www.health.nsw.gov.au/
    \item Centers for Disease Control and Prevention. \textit{Haemophilus influenzae disease.} https://www.cdc.gov/hi-disease/
    \item NHS. \textit{Epiglottitis.} https://www.nhs.uk/conditions/epiglottitis/
    \item MSD Manual. \textit{Epiglottitis: professional overview.} https://www.msdmanuals.com/
\end{itemize}
